\section*{Capítulo \ref{ch:g7:how-animals-breathe} --- Cómo respiran los animales}

\begin{solution}{exo:g7:how-animals-breathe:1}
Una superficie de intercambio grande, fina y húmeda donde la \omterm{prop:g4:journey-of-food:absorb}{sangre} (o
el líquido corporal) se encuentra con el \omterm{def:g6:exploring-our-environment:components}{entorno}. Grande: paso total
suficiente; fina: paso fácil; húmeda: los gases pasan disueltos.
\end{solution}

\begin{solution}{exo:g7:how-animals-breathe:2}
\omterm{def:g7:how-animals-breathe:four}{Pulmones}: \omterm{prop:g3:sorting-living-things:groups}{mamíferos}, \omterm{prop:g3:sorting-living-things:groups}{aves} (también \omterm{prop:g4:classifying-animals:classes}{reptiles} y \omterm{prop:g3:sorting-living-things:groups}{anfibios} adultos).
\omterm{def:g7:how-animals-breathe:four}{Branquias}: \omterm{prop:g3:sorting-living-things:groups}{peces}, \omterm{ex:g2:animals-grow-up:frog}{renacuajos} (también mejillones y larvas acuáticas).
\omterm{def:g7:how-animals-breathe:four}{Tráqueas}: saltamontes, abejas --- los \omterm{prop:g3:sorting-living-things:groups}{insectos} en general. \omterm{def:g7:how-animals-breathe:four}{Respiración
cutánea}: la lombriz de tierra; y los \omterm{prop:g3:sorting-living-things:groups}{anfibios} como complemento.
\end{solution}

\begin{solution}{exo:g7:how-animals-breathe:3}
El agua entra por la boca, se empuja sobre las \omterm{def:g7:how-animals-breathe:four}{branquias} plumosas ---
donde ocurre el intercambio --- y sale por debajo de los opérculos: una
corriente de un solo sentido, bombeada por la boca y los opérculos
alternándose.
\end{solution}

\begin{solution}{exo:g7:how-animals-breathe:4}
Por las \omterm{def:g7:how-animals-breathe:four}{tráqueas}: una red de tubos de aire finos que van desde unas
aberturas repartidas por el cuerpo y se ramifican hasta todos los
órganos --- el aire se reparte puerta a puerta, así que la \omterm{prop:g4:journey-of-food:absorb}{sangre} tiene
poco que llevar.
\end{solution}

\begin{solution}{exo:g7:how-animals-breathe:5}
Fuera del agua, sus superficies branquiales plumosas se aplastan y se
pegan unas a otras: la gran superficie se reduce a casi nada, y sin
superficie no hay paso --- \omterm{def:g7:what-is-respiration:respiration}{oxígeno} de sobra y ninguna máquina para
tomarlo.
\end{solution}

\begin{solution}{exo:g7:how-animals-breathe:6}
La \omterm{def:g7:how-animals-breathe:four}{respiración cutánea} --- y falla la condición de \omterm{def:g6:exploring-our-environment:conditions}{humedad}: la \omterm{def:g1:the-five-senses:sense}{piel} que
se seca ya no puede dejar pasar gases y la lombriz se asfixia al aire
libre. Ser grande y fina no salva a una superficie que ya no está
húmeda.
\end{solution}

\begin{solution}{exo:g7:how-animals-breathe:7}
Renacuajo: \omterm{def:g7:how-animals-breathe:four}{branquias}. Adulta: \omterm{def:g7:how-animals-breathe:four}{pulmones} sencillos más la \omterm{def:g1:the-five-senses:sense}{piel} húmeda ---
y la \omterm{def:g1:the-five-senses:sense}{piel} sola sirviendo durante el invierno amodorrado en el fondo de
la charca.
\end{solution}

\begin{solution}{exo:g7:how-animals-breathe:8}
Su máquina es el \omterm{def:g7:how-animals-breathe:four}{pulmón} --- heredado de antepasados terrestres que
alimentaban con \omterm{def:g2:animals-grow-up:born}{leche} --- así que tiene que salir a la superficie a
respirar aire, por muy de \omterm{prop:g3:sorting-living-things:groups}{pez} que sea su vida: si se lo mantiene
debajo, se ahoga. El vestido es del agua; la máquina, de la familia.
\end{solution}

\begin{solution}{exo:g7:how-animals-breathe:9}
Una reserva de aire que lleva encima --- una botella de buceo bajo los
élitros, de la que tiran sus \omterm{def:g7:how-animals-breathe:four}{tráqueas}. Tiene que volver a la superficie
periódicamente para renovar la burbuja.
\end{solution}

\begin{solution}{exo:g7:how-animals-breathe:10}
(a) Mejillón: residente del agua, sin movimientos visibles pero con
\omterm{def:g7:how-animals-breathe:four}{branquias} dentro de la \omterm{def:g1:animals-around-us:covering}{concha}, lavadas por una corriente bombeada. (b)
Saltamontes: residente del aire; filas de aberturas diminutas, abdomen
que late --- \omterm{def:g7:how-animals-breathe:four}{tráqueas}. (c) Culebra: un reptil que va y viene ---
\omterm{def:g7:how-animals-breathe:four}{pulmones}; nada con los orificios nasales fuera y sale a la superficie a
respirar.
\end{solution}

\begin{solution}{exo:g7:how-animals-breathe:11}
Las \omterm{def:g7:how-animals-breathe:four}{tráqueas} reparten con tubos de aire ramificados, magníficos en
milímetros y demasiado lentos en distancias largas: en un cuerpo del
tamaño de una ballena, las entrañas quedarían demasiado lejos de
cualquier abertura. La máquina pone un techo de tamaño, y los \omterm{prop:g3:sorting-living-things:groups}{insectos}
viven por debajo de él.
\end{solution}

\begin{solution}{exo:g7:how-animals-breathe:12}
El pliego: una superficie de encuentro grande, fina y húmeda entre el
cuerpo y el \omterm{def:g6:exploring-our-environment:components}{entorno}. Las \omterm{def:g7:how-animals-breathe:four}{branquias} lo cumplen con superficies plumosas
desplegadas por el agua; los \omterm{def:g7:how-animals-breathe:four}{pulmones}, con millones de bolsas plegadas
a resguardo y en húmedo; las \omterm{def:g7:how-animals-breathe:four}{tráqueas}, con la superficie interior sumada
de sus puntas más finas, húmedas en los extremos; y la \omterm{def:g1:the-five-senses:sense}{piel}, con todo el
exterior húmedo del cuerpo mantenido grande respecto de un cuerpo
pequeño y delgado. Cuatro soportes, una sola cosa sostenida.
\end{solution}

\begin{solution}{pb:g7:how-animals-breathe:1}
\textbf{1.} Alevines de trucha y espinosos: \omterm{def:g7:how-animals-breathe:four}{branquias}. \omterm{ex:g2:animals-grow-up:frog}{Renacuajos}:
\omterm{def:g7:how-animals-breathe:four}{branquias} (los que tienen patas, a media reconstrucción). Ranas
adultas: \omterm{def:g7:how-animals-breathe:four}{pulmones} más \omterm{def:g1:the-five-senses:sense}{piel}. Caracoles de charca: \omterm{def:g7:how-animals-breathe:four}{pulmones} --- van y
vienen a la superficie. Escarabajo buceador: \omterm{def:g7:how-animals-breathe:four}{tráqueas}, alimentadas por
la burbuja que lleva encima. Larvas de efímera: \omterm{def:g7:how-animals-breathe:four}{branquias} (los penachos
plumosos). Lombrices: \omterm{def:g1:the-five-senses:sense}{piel}.

\textbf{2.} \omterm{def:g7:how-animals-breathe:four}{Branquias}. La agitación renueva el agua contra la
superficie branquial --- una corriente fabricada por ella misma donde
la charca quieta no da ninguna, para mantener en el punto de paso agua
fresca cargada de \omterm{def:g7:what-is-respiration:respiration}{oxígeno}.

\textbf{3.} Los caracoles de charca, el escarabajo buceador, las ranas
adultas (y cualquier culebra de visita): todos los que usan \omterm{def:g7:how-animals-breathe:four}{pulmones} o
burbuja --- sus máquinas toman el \omterm{def:g7:what-is-respiration:respiration}{oxígeno} del aire, así que la
superficie es su gasolinera.

\textbf{4.} La reconstrucción respiratoria de la \omterm{def:g2:animals-grow-up:metamorphosis}{metamorfosis}: las
\omterm{def:g7:how-animals-breathe:four}{branquias} que servían al renacuajo sin patas se están cambiando por los
\omterm{def:g7:how-animals-breathe:four}{pulmones} (y la \omterm{def:g1:the-five-senses:sense}{piel} respiratoria) que necesitará la ranita de tierra.

\textbf{5.} Lo pasan mal primero los de \omterm{def:g7:how-animals-breathe:four}{branquias}: todo su suministro
es la reserva disuelta del agua, que va menguando. Los que van y vienen
--- caracol, escarabajo, rana --- tiran del aire ilimitado de arriba y
solo tienen que subir más a menudo.

\textbf{6.} El milímetro de arriba toca el aire: allí el \omterm{def:g7:what-is-respiration:respiration}{oxígeno} se
disuelve desde arriba, así que la capa de la película es el agua menos
pobre de una charca que se asfixia. Los espinosos están haciendo cola
en el último pozo.

\textbf{7.} Se atascan las \omterm{def:g7:how-animals-breathe:four}{branquias}: el barro fino se posa en las
superficies plumosas y bloquea el paso. Es de esperar que los que
respiran por \omterm{def:g7:how-animals-breathe:four}{branquias} huyan de la nube y que el bombeo de opérculos se
vuelva trabajoso --- con espasmos de tos para despejar el limo.

\textbf{8.} La \omterm{def:g7:how-animals-breathe:four}{respiración cutánea}, y la condición de húmedo pero sin
aire: las galerías encharcadas tienen demasiado poco \omterm{def:g7:what-is-respiration:respiration}{oxígeno}, así que
las lombrices salen a respirar --- cambian asfixiarse abajo por el
riesgo de secarse arriba.

\textbf{9.} La tabla, tal como se pide (máquina y medio por \omterm{def:g1:animals-around-us:animal}{animal}; van
y vienen: caracol, escarabajo, rana). Encabezando todas las columnas:
la superficie de intercambio grande, fina y húmeda --- el único pliego
de \omterm{def:g6:exploring-our-environment:conditions}{condiciones} que sostienen las cuatro máquinas.

\textbf{10.} Ranas: cubiertas --- amodorradas en el barro y respirando
por la \omterm{def:g1:the-five-senses:sense}{piel} su pequeña necesidad invernal. Caracoles y escarabajo: con
problemas --- su superficie está sellada; pasan el invierno
aletargados, con las necesidades bajadas o con burbujas atrapadas.
\omterm{prop:g3:sorting-living-things:groups}{Peces}: cubiertos mientras dure la reserva del agua --- el frío frena
todos los ruidos de motor (\cref{ex:g7:what-is-respiration:demand}), y
las \omterm{def:g7:how-animals-breathe:four}{branquias} sirven a la demanda reducida bajo el hielo.

\textbf{11.} La campana de seda de la araña es un depósito de aire fijo
donde la burbuja del escarabajo es uno portátil: los dos guardan aire
de la superficie bajo el agua y los dos siguen sirviendo a una máquina
de aire (el equipo de \omterm{def:g7:how-animals-breathe:four}{pulmones} en libro y \omterm{def:g7:how-animals-breathe:four}{tráqueas} de la araña) en el
medio equivocado.

\textbf{12.} Porque la charca ofrece \omterm{def:g7:what-is-respiration:respiration}{oxígeno} en dos medios --- escaso y
disuelto abajo, abundante arriba --- y cada máquina cosecha un medio a
un tamaño de vida: solo con las cuatro trabajando queda habitado cada
rincón, del barro a la película de la superficie.
\end{solution}
